\documentclass{report}
\usepackage{amsfonts, amsmath}
\newcommand\R{{\mathbb R}}
\newcommand\Q{{\mathbb Q}}
\newcommand\N{{\mathbb N}}
\newcommand\Or{{\mathcal O}}
\newcommand\ve{\varepsilon}
%\pagestyle{empty}
\begin{document}
\large
\centerline{\Huge Analisi Matematica II}
\renewcommand{\thefootnote}{ } 
\renewcommand{\thefootnote}{\arabic{footnote}} 
\vskip.1cm
\centerline{\Large Soluzioni esonero 23-04-2003}
\vskip1cm
\begin{enumerate}
\item La solita figura che confronta l'area sotto la curva $x^{-\frac
12}$ e la somma permette di stabilire
\[
\int_1^{N+1}\frac 1{\sqrt x} dx\leq \sum_{n=1}^N\frac 1{\sqrt n}\leq
1+ \int_1^{N}\frac 1{\sqrt x} dx
\]
Dunque
\[
\sum_{n=1}^N\frac 1{\sqrt n}\geq  2\sqrt{N+1}-2 >1000
\]
se $N\geq 501^2-1$. Abbiamo cosi un $N$ per cui la somma supera mille.

Se vogliamo avere una stima del primo intero $\bar N$ per cui la somma
supera mille allora notiamo che quanto gia detto implica $\bar N\leq
501^2-1$, inoltre
\[
1000\leq \sum_{n=1}^{\bar N} \frac 1{\sqrt n}\leq 1+2\sqrt{\bar N}-2,
\]
dunque $\bar N\geq (500.5)^2$. Questo significa che $\bar N$ \`e ora
noto a meno di circa $500$. 

Quanto sopra \`e pi\`u che sufficiente come soluzione dell'esercizio.

Comunque, per i curiosi, vediamo come si pu\`o ottenre una precisione
assai maggiore.

Sia definisca
\[
\Delta(N,L):=\sum_{n=N}^L \frac 1{\sqrt n}-\int_N^{L+1}\frac 1{\sqrt x}
dx.
\]
Come vedremo tra poco esiste 
\[
\Delta(N,\infty):=\lim_{L\to\infty}\Delta(N,L).
\]
La rilevanza di tale quantit\`a risiede nella ovvia formula
\[
\Delta (1,\infty)=\Delta(1,N)+\Delta(N+1,\infty)
\]
da cui
\begin{equation}\label{eq:f1}
\sum_{n=1}^N \frac 1{\sqrt n}=\int_1^{N+1}\frac 1{\sqrt
x}dx +\Delta(1,\infty)-\Delta(N+1,\infty).
\end{equation}
Per una buona stima della somma occorre dunque una buona stima di
$\Delta$.
\begin{align*}
&\Delta(N,L)=\sum_{n=N}^L\int_n^{n+1}\left[\frac 1{\sqrt n}-\frac
1{\sqrt x}\right]  dx\\
&=\sum_{n=N}^L\int_n^{n+1}\left[ \frac 12
n^{-\frac 32}(x-n)-\frac 38n^{-\frac 52}(x-n)^2+\frac 5{16}\xi^{-\frac
72}(x-n)^3\right] dx\\
&=\sum_{n=N}^L \left[\frac 14 n^{-\frac 32}-\frac 18n^{-\frac 52}+{\mathcal
O}(\frac 5{64}n^{-\frac 72})\right]
\end{align*}
dove per ${\mathcal O}(\ve)$ si intende un quantit\`a minore di $\ve$ e
maggiore di $-\ve$. Come annunciato il limite per $L\to\infty$ \`e ben
definito (le serie convergono). 
Per continuare \`e necessario stimare le tre serie di cui sopra. 

Visto che sono tutte dello stesso tipo facciamolo una vota per tutte.
Sia $p> 1$, allora
\[
\sum_{n=N}^\infty n^{-p}\leq \int_{_1}^\infty
x^{-p}dx=\frac{(N-1)^{1-p}}{p-1}.
\]
Ma si puo fare meglio
\begin{align*}
\sum_{n=N}^\infty n^{-p}&=\int_N^\infty
x^{-p}dx+\sum_{n=N}^\infty\int_n^{n+1}(n^{-p}-x^{-p}) dx\\
&=\frac{N^{1-p}}{p-1}+\sum_{n=N}^\infty{\mathcal O}(\frac p2 n^{-p-1})
=\frac{N^{1-p}}{p-1}+{\mathcal O}\left(\frac 1{2(N-1)^p}\right).
\end{align*}
E ancora meglio
\begin{align*}
\sum_{n=N}^\infty n^{-p}
&=\frac{N^{1-p}}{p-1}+\sum_{N}^\infty \int_n^{n+1}\left[
pn^{-p-1}(x-n)-\frac{p(p+1)}2 \xi^{-p-2}(x-n)^2] dx\right]\\
&=\frac{N^{1-p}}{p-1}+\frac p2\sum_{N}^\infty
n^{-(p+1)}+\frac{p(p+1)}4{\mathcal O}\left(\sum_{n=N}^\infty n^{-(p+2)}\right)\\
&=\frac{N^{1-p}}{p-1}+\frac p2\left[\frac{N^{-p}}p+\Or\left(\frac
1{2(N_1)^{p+1}}\right)\right] +\Or\left(\frac p{4(N-1)^{p+1}}\right)\\
&=\frac{N^{1-p}}{p-1}+\frac 1{2N^p}+\Or\left(\frac{p+2}{4(N-1)^{p+1}}\right).
\end{align*}
Ovviamente, si pu\`o fare ancora meglio usando uno sviluppo di taylor
ad un ordine maggiore, purtuttavia quanto sopro \`e sufficiente per i
nostri scopi per cui ci fermiamo qui.

Dunque
\begin{equation}\label{eq:f2}
\Delta(N,\infty)=\frac 1{2\sqrt N}+\frac 1{24 N^{\frac
32}}+\Or\left(\frac {27}{128 (N-1)^{\frac 52}}\right)
\end{equation}
per cui
\[
\Delta(16,\infty)=\frac 1{2\sqrt 16}+\frac 1{24\times 16^{\frac 32}}+{\mathcal
O}(2.5 \times 10^{-4})
\]
e, sommando a mano i primi quindici termini, si ottiene
\begin{equation}\label{eq:f3}
\Delta(1,\infty)=0.53965+{\mathcal O}(2.5 \times 10^{-4}).
\end{equation}
Finalmente, collezionando le uguaglianze \eqref{eq:f1}, \eqref{eq:f2},
\eqref{eq:f3}, possiamo scrivere 
\begin{align*}
\sum_{n=1}^N \frac 1{\sqrt n}=&2\sqrt{N+1}-2
+0.53965-\frac1{2\sqrt{N+1}}-\frac 1{24 (N+1)^{\frac 32}}\\
&\ +{\mathcal
O}(\frac{27}{128 N^{\frac 52}}+2.5 \times 10^{-4}) .
\end{align*}
Dunque si ha
\begin{align*}
\sum_{n=1}^{250730} \frac 1{\sqrt n}&=000.99958+{\mathcal O}(2.6\times 10^{-4})\\
\sum_{n=1}^{250731} \frac 1{\sqrt n}&=1000.0016+{\mathcal O}(2.6\times 10^{-4}) .
\end{align*}
Si pu\`o quindi affermare che l'intero per cui la somma \`e, per la prima volta, pi\`u
grande di 1000 \`e $\bar N=250731$.\footnote{Se volete una verifica
indipendente, potete scrivere un
piccolo programma (nel vostro linguaggio di programmazione preferito)
che trova $\bar N$ sommando e fermandosi quando la 
somma \`e maggiore di $1000$. Notate comunque che se invece di mille
il limite da superare fosse $10^9$ allora $\bar N$ sarebbe dell'ordine
di grandezza di $10^{18}$ e difficilmente un calcolatore potrebbe fare
$10^{18}$ estrazioni di radici quadrate, inversioni e somme nell'arco
della vostra vita, mentre una semplice sofisticazione del metodo sopra
descritto vi darebbe la risposta con poche ore di calcolo.}

\item L'inverso  raggio di convergenza \`e dato dalla formula
\[
\lim_{n\to \infty}[e^{\sqrt n}]^{\frac 1n}=\lim_{n\to \infty} e^{\frac
1{\sqrt n}}=1.
\]
In alternativa
\[
\lim_{n\to\infty} \frac{e^{\sqrt{n+1}}}{e^{\sqrt n}}=\lim_{n\to\infty}
e^{\sqrt{n+1}-\sqrt n} =\lim_{n\to\infty}e^{\frac1{\sqrt{n+1}+\sqrt
n}}=1.
\]
\item Si calcolino i coefficienti di Fourier
\[
a_0=\frac 1\pi\int_{-\pi}^\pi \cos\frac x2 dx=\frac 2\pi\int_{-\frac
\pi 2}^{\frac \pi 2}\cos x dx=\frac 4\pi.
\]
Tutti is $b_n$ sono zero per parit\`a mentre gli $a_n$ si ottengono
integrando due volte per parti (come al solito) ottenendo
\[
a_n=\frac {(-1)^n 4}{\pi (1-4n^2)}.
\]
Poche la funzione $g$ \`e peiodica di periodo $2\pi$, per costruzione,
e coincide con $\cos x/2$ nell'intervallo $[-\pi,\pi]$ segue che
$g(x)=|\cos x/2|$ per ogni $x\in \R$. Dunque $g$ \`e continua ma non
ovunque derivabile.
\item 
\[
\lim_{L\to\infty}\frac 12\int_2^L\frac 1{x-1}-\frac 1{x+1} dx=
\frac 12 \lim_{L\to\infty}\left \{\ln\frac{L-1}{L+1}-\ln\frac
13\right\}
=\ln\sqrt 3.
\]
\end{enumerate}
\end{document}
